\uuid{GHUY}
\titre{Points stationnaires}
\auteur{hanani ; mzari}
\organisation{exo7}
\licence{CC BY-NC-SA 4.0}
\datecreate{2020-01-20}
\niveau{L1}
\module{Géométrie}
\chapitre{Courbes planes}
\sousChapitre{Courbes paramétrées}
\difficulte{3}
\theme{Courbes paramétrées, points stationnaires}
\exo7qcmid{lille-2.courbes.010}

\contenu{
\begin{question}
\qclassification{244.01}
\qtags{motcle=points stationnaires, asymptotes}
La trajectoire $\Gamma$ d'une particule en mouvement est donnée par les équations
$$x=\frac{t^2}{1-t^2}\quad \mbox{et} \quad  y=\frac{t^2}{1+t}.$$
\begin{answers}
\bad{Le point de paramètre $0$ est un point de rebroussement de seconde espèce.}
\good{La tangente à $\Gamma$ au point de paramètre $1$ est la droite d'équation $y=x$.}
\bad{$\Gamma$ admet la droite d'équation $y=2$ comme asymptote quand $t$ tend vers $1$.}
\good{$\Gamma$ admet la droite d'équation $y=2x-1/2$ comme asymptote quand $t$ tend vers $-1$.}
\end{answers}
\begin{explanations}
On a : $f(t)=t^2(1,1)+t^3(0,-1)+t^3\left(\varepsilon_1(t),\varepsilon_2(t)\right)$ avec $\displaystyle \lim _{t\to 0}\varepsilon_i(t)=0$. Donc le point de paramètre $0$ est un point de rebroussement de première espèce et la tangente en ce point est la droite d'équation $y=x$. Enfin, l'étude des branches infinies, quand $t$ tend vers $1$ et quand $t$ tend vers $-1$, montre que les droites d'équation $y=1/2$ et d'équation $y=2x-1/2$ sont des asymptotes.
\end{explanations}
\end{question}
}
